%% =============================================================================
%% SUBSTITUIR o ambiente \begin{figure}...\end{figure} de
%% fig:metodologia-visao-geral (Cap. 4, Seção "Visão Geral da Abordagem")
%% por este. Engloba as DUAS campanhas, na paleta nude. O \ref no texto
%% continua válido (mesmo label). Requer as cores palPaper/palFill/palStrong/
%% palEdge/palEdgeDark (definidas em 000-main.tex) e as bibliotecas
%% positioning, shapes.geometric, arrows.meta (já carregadas).
%% =============================================================================
\begin{figure}[H]
\centering
\resizebox{\linewidth}{!}{%
\begin{tikzpicture}[
  font=\linespread{1}\footnotesize,
  >={Stealth[length=2mm]},
  node distance=4mm,
  base/.style={draw=palEdge, rounded corners=2pt, align=center, inner sep=3pt,
    fill=palPaper},
  etapa/.style={base, text width=5.8cm, minimum height=7mm},
  firm/.style={etapa, fill=palFill},
  nova/.style={base, text width=4.8cm, minimum height=7mm, draw=palEdgeDark,
    fill=palStrong},
  dir/.style={base, text width=4.8cm, minimum height=7mm},
  conf/.style={base, text width=6.2cm, fill=palFill, draw=palEdgeDark, very thick},
  dec/.style={draw=palEdge, diamond, aspect=2.4, inner sep=1pt, fill=palPaper,
    align=center},
  out/.style={base, text width=3.0cm, minimum height=10mm, fill=palStrong,
    draw=palEdgeDark},
  outn/.style={base, text width=3.0cm, minimum height=10mm, fill=palPaper},
  rot/.style={font=\scriptsize\itshape, align=center}]

  % --- espinha (i)-(v)
  \node[etapa] (e1) {(i) Seleção de módulos críticos};
  \node[firm, below=of e1] (e2) {(ii) Instrumentação com \textit{stubs} de \textit{hardware}};
  \node[firm, below=of e2] (e3) {(iii) Especificação de propriedades\\(\textit{datasheet}, documentação, código)};
  \node[etapa, below=of e3] (e4) {(iv) Configuração e execução do ESBMC};
  \node[etapa, below=of e4] (e5) {(v) Análise de contraexemplos};

  % --- conferencia compartilhada, com os dois modos
  \node[conf, below=6mm of e5] (conf) {\textbf{Conferência contra o código de produção}\\[1pt]
    1.\textsuperscript{a} campanha: leitura do código\\
    2.\textsuperscript{a} campanha: \textit{replay} do código real {\footnotesize+} ASan/UBSan};
  \node[dec, below=6mm of conf] (dec) {Reproduz em\\produção?};

  % --- desfechos
  \node[out, below=10mm of dec] (def) {Defeito reproduzido / VAM};
  \node[outn, left=20mm of def] (fp) {Falso positivo (modelagem / ferramenta)};
  \node[outn, right=20mm of def] (ok) {Sem violação até o limite $k$};

  % --- coluna direita (vi)-(viii)
  \node[dir, right=2.4cm of e1.north east, anchor=north west] (e6) {(vi) Reporte à comunidade PX4\\(condicionado à revisão)};
  \node[nova, below=of e6] (e7) {(vii) Validação cruzada (CBMC)\\e verificação de \textit{patch}};
  \node[nova, below=of e7] (e8) {(viii) Reprodução (ArduPilot)\\e integração contínua};

  % --- setas da espinha e da coluna direita
  \foreach \a/\b in {e1/e2,e2/e3,e3/e4,e4/e5,e5/conf,conf/dec,e6/e7,e7/e8}
    \draw[->] (\a) -- (\b);

  % --- decisao -> desfechos
  \draw[->] (dec) -- node[rot, right] {sim} (def);
  \draw[->] (dec.west) -| node[rot, pos=0.25, above] {guarda\\omitida} (fp.north);
  \draw[->] (dec.east) -| node[rot, pos=0.25, above] {nenhuma\\violação} (ok.north);

  % --- defeito/VAM -> reporte (vi)
  \draw[->] (def.east) .. controls +(2.6,0) and +(0,-2.2) .. (e6.south);

  % --- legenda
  \node[firm, text width=3mm, minimum height=3mm, below=12mm of fp.south west,
    anchor=north west] (l1) {};
  \node[right=1mm of l1, font=\scriptsize, anchor=west] (t1) {etapas de \textit{firmware}};
  \node[nova, text width=3mm, minimum height=3mm, right=5mm of t1] (l2) {};
  \node[right=1mm of l2, font=\scriptsize, anchor=west] (t2) {acréscimos deste trabalho};
  \node[conf, text width=3mm, minimum height=3mm, right=5mm of t2] (l3) {};
  \node[right=1mm of l3, font=\scriptsize, anchor=west] {conferência (1.\textsuperscript{a} e 2.\textsuperscript{a} campanhas)};
\end{tikzpicture}}
\caption{Metodologia de verificação em oito etapas, aplicada em duas campanhas.
As etapas~(i) e~(iv)--(vi) seguem o fluxo de Sousa et al.~\cite{Sousa2026}; as
etapas~(ii) e~(iii) são específicas de \textit{firmware} e as etapas~(vii)
e~(viii) são acréscimos deste trabalho. A conferência contra o código de
produção (etapa~v) é feita por leitura do código na primeira campanha e por
\textit{replay} do código real sob sanitizadores (ASan/UBSan) na segunda. Apenas
defeitos reproduzidos e violações alcançáveis no modelo (VAM) seguem para o
reporte~(vi), que depende de revisão; falsos positivos e resultados sem violação
encerram o fluxo do módulo. Os desfechos usam a escala do
Capítulo~\ref{cap:introducao}. Elaborado pela autora.}
\label{fig:metodologia-visao-geral}
\end{figure}
